%% ============================================================
%%  CHAPITRE 5 : IMPLÉMENTATION ET RÉALISATION
%% ============================================================
\chapter{Implémentation et Réalisation}
\label{chap:implementation}

Ce chapitre détaille les aspects techniques de l'implémentation de la plateforme MASVS Audit Copilot. Nous y aborderons la sécurité de l'API, le moteur d'intégration de l'IA, et la logique de génération des rapports d'audit.

\section{Sécurisation de l'API (FastAPI)}
La sécurité de l'outil d'audit lui-même est primordiale. L'API FastAPI est protégée par un système d'authentification basé sur des \textbf{JSON Web Tokens (JWT)}.

Lors de la connexion, l'utilisateur s'authentifie via le protocole OAuth2. L'API génère un token JWT signé avec un algorithme cryptographique asymétrique (ou symétrique fort selon la configuration) et une durée de vie limitée. Ce token doit être inclus dans l'en-tête \code{Authorization: Bearer <token>} de chaque requête subséquente. De plus, les mots de passe des utilisateurs sont hachés en base de données à l'aide de l'algorithme \textbf{Bcrypt}, garantissant qu'ils ne puissent pas être récupérés en clair, même en cas de compromission de la base de données.

\section{Implémentation du Triage IA (Module LLM Client)}
Le cœur de la réduction des faux positifs réside dans l'interaction avec le modèle de langage via Ollama. L'implémentation a été conçue pour être robuste et "fail-safe".

Le module \code{LLMClient} gère l'orchestration des requêtes vers l'IA. Pour chaque vulnérabilité détectée par MobSF, le client génère un "prompt système" strict contraignant le LLM à répondre uniquement sous forme d'un objet JSON. Ce prompt inclut le contexte complet : titre de la faille, description, catégorie MASVS et extrait de code source.

\begin{tcolorbox}[codebox, title={Extrait : Logique de Fallback (llm\_triage.py)}]
\begin{minted}[fontsize=\footnotesize]{python}
class LLMClient:
    def __init__(self):
        self.provider = settings.LLM_PROVIDER # "ollama" ou "openai"
        self.ollama_url = settings.OLLAMA_URL

    async def analyze_finding(self, finding: Finding) -> dict:
        prompt = self._build_prompt(finding)
        try:
            if self.provider == "openai":
                return await self._call_openai(prompt)
            else:
                return await self._call_ollama(prompt)
        except Exception as e:
            # Mécanisme de fallback : Si l'API cloud échoue (ex: erreur 401),
            # le système bascule automatiquement sur Ollama local.
            logger.warning(f"Provider {self.provider} failed: {e}. Fallback to Ollama.")
            return await self._call_ollama(prompt)
\end{minted}
\end{tcolorbox}

Le JSON retourné par l'IA est ensuite parsé de manière sécurisée. S'il indique que la faille est un faux positif (\code{"is\_true\_positive": false}), la sévérité de la vulnérabilité est automatiquement dégradée à "INFO" et elle est exclue du calcul du score MASVS.

\section{Module d'Exportation des Rapports}
Une fois le pipeline de scan terminé (MobSF $\rightarrow$ Mapping $\rightarrow$ SBOM $\rightarrow$ IA), le système génère les livrables finaux. Le projet intègre un module d'export modulaire, capable de générer trois formats de rapports :

\begin{itemize}[label=\textcolor{PrimaryBlue}{\faFileExport}]
    \item \textbf{Markdown (\code{markdown\_exporter.py}) :} Génère un fichier texte brut formaté, très utile pour l'intégration dans des wikis internes ou des \textit{issues} GitHub/GitLab.
    \item \textbf{PDF (\code{pdf\_generator.py}) :} Utilise la bibliothèque \code{ReportLab} pour générer un document PDF professionnel prêt à être envoyé aux parties prenantes ou à la direction. Il inclut un sommaire, des tableaux de synthèse codés par couleur et le détail complet de chaque vulnérabilité avec sa remédiation.
    \item \textbf{SARIF 2.1 :} Le \textit{Static Analysis Results Interchange Format} est un standard JSON pour les outils d'analyse statique. Cet export permet d'ingérer directement les résultats dans des plateformes comme GitHub Advanced Security ou SonarQube, intégrant ainsi pleinement l'outil dans une chaîne CI/CD (DevSecOps).
\end{itemize}

\section{Interface en Ligne de Commande (CLI)}
Pour les ingénieurs en cybersécurité et l'automatisation, une interface graphique n'est pas toujours optimale. Nous avons développé une CLI complète en Python en utilisant la bibliothèque \textbf{Typer}.

Cette CLI permet d'interagir avec l'API REST de manière scriptable. Elle supporte l'authentification, le déclenchement de scans, le suivi de la progression et le téléchargement des rapports. L'interface CLI gère l'affichage asynchrone via des barres de progression interactives (\code{rich.progress}), offrant un retour visuel direct dans le terminal de l'auditeur.

\cleardoublepage
